\documentclass[12pt]{article}
\usepackage[utf8]{inputenc}
\usepackage[T1]{fontenc}
\usepackage{hyperref}
\title{Miarr\~a Language Resources for K-12 Curriculum Development}
\author{Compiled by Claude Sonnet 5.5 (claude-sonnet-5-5) for Andr\'as Kornai}
\date{}
\begin{document}
\maketitle

\section{Introduction}
Miarr\~a (ISO 639-3 xmi; Glottocode miar1236) is an unattested language of South America. Search found almost nothing. The only item mentioning it is a Glottolog bibliographic record for an Orlando Villas Boas report (1975) on groups in the Xingu area. It says, citing the Kuikuru, that Miarr\~a speakers use the same language as the Yawalapiti. This is not to be confused with Himarim\~a (hir), a different language.

\section{Primary Educational Resources}
None exist as far as could be found. No speakers, texts or word lists were found.

\section{Linguistic and Reference Materials}
\begin{itemize}
\item Glottolog reference record, Villas Boas 1975 (mentions xmi): \url{https://glottolog.org/resource/reference/id/120022}
\end{itemize}

\section{Religious and Community Resources}
None found.

\section{Digital and Audio Resources}
None found.

\section{Curriculum Development Context}
A K-12 curriculum is not feasible for an unattested language with no documented speakers or data. The most that can be said is that the Glottolog record points to a possible link with Yawalapiti, which would need verification from specialist sources.
\end{document}